%% CPP 2027 Submission — An Empirical Trail to Formalized Pat Theory
%% Compile: pdflatex + bibtex (acmart)
\documentclass[sigplan,10pt,anonymous,review]{acmart}
\settopmatter{printfolios=true,printacmref=false}

\AtBeginDocument{%
  \providecommand\BibTeX{{\normalfont B\scshape ib\TeX}}}

\settopmatter{printacmref=false}
\renewcommand\footnotetextcopyrightpermission[1]{}
\pagestyle{plain}

\usepackage[utf8]{inputenc}
\usepackage[T1]{fontenc}
\usepackage{amsmath}
\ifdefined\XeTeXinterchartokenstate
  % XeLaTeX: unicode-math (loaded by acmart) supplies the amssymb glyphs
\else
  \usepackage{amssymb}
\fi
\usepackage{booktabs}
\usepackage{microtype}

\begin{document}

\title{An Empirical Trail to Formalized Pat Theory:\\
from 0/1 OOD Observations to Machine-Checked Theorems\\
(complete research path, timestamped evidence, all PROVED, no \texttt{sorry})}

\author{Anonymous}
\affiliation{%
  \institution{Independent}
  \country{}}


\date{2026-08-20}

\begin{abstract}
We record a complete research trail --- from a transformer training observation
(2026-07-29) through a systematic experiment series (E1--E18), the extraction of
structural laws, the construction of the Pat framework (basepoint-relative
operator semantics), and its full Lean 4 + mathlib formalization (claims
R136--R153 and beyond; all PROVED, no \texttt{sorry}; full \texttt{lake build}
passes, 3631 jobs). The trail is anchored by timestamped evidence: session
records, a 2303-turn task log, git history, and experiment artifacts. The
starting observation is the 0/1 OOD dichotomy: a low-epoch (60) transformer
reaches training accuracy 1.000 yet generalizes 0/32 on a same-basepoint
notation change, and 32/32 after equivalence-declaration bridging --- a
reproducible two-pole pattern across 45 training runs. From this observation
the trail extracts five structural laws (anchor binding, notation
translatability, shift invariance, structure-bearing representation,
definition-layer invisibility), maps them to framework claims (R136 direction
declaration, R138 phase locking, R143 interlock matrix), and formalizes the
chain to the continuum construction (R150, R151, R163). We contribute (i) the
complete experiment-to-formalization path as a reproducible artifact --- each
experiment documented individually with design reason, procedure, result, and
significance --- and (ii) a methodology for recording research trails under
provenance pressure: SHA-256-anchored manifests, per-claim publication, and a
self-contained reproducibility package.
\end{abstract}

\keywords{Lean, formalization, transformer, OOD, basepoint, reproducibility, research trail}

\maketitle

\section{Introduction}

This paper is a complete, machine-checkable record of a research trail: how a
transformer training observation became a set of structural laws, then a
framework, then Lean-formalized theorems. We record the trail because it is
itself the contribution: the path from an empirical two-pole OOD pattern to a
formalized theory of basepoint-relative operator semantics, with every step
timestamped and reproducible.

The trail starts on 2026-07-29 with a first \emph{imperfect} generalization
observation (relation-leave-one-out generalization, 22:39:12 UTC+8) ---
imperfect, yet already sufficient to guide the exploration direction. It
reaches its first \emph{perfect} generalization on 2026-08-15 04:01 (all 19
operations OOD 0.000 $\to$ 1.000). The core two-pole pattern (E12) is: a
60-epoch, 2-layer transformer (dim 64) reaches training accuracy 1.000, yet
scores \textbf{0/32} on a same-basepoint notation change (succ$_a \to$
succ$_b$) and \textbf{32/32} after adding equivalence-declaration bridging
samples --- reproduced across 45 training runs (5 groups $\times$ 3 seeds
$\times$ 3 independent runs), with no deviation.

From this pattern, a systematic experiment series (E1--E18) extracts five
structural laws; the laws map to the Pat framework (basepoint-relative
operators: every operator has its own basepoint, directions must be declared
in pairs); the framework is formalized in Lean 4 (mathlib v4.32.2), claims
R136--R153 all PROVED with no \texttt{sorry} (completed 2026-08-13 13:31), and
the chain extends to the continuum construction (R150 Wang's phase-locking
theorem, R151 continuum as Pat-grid closure, R163 endogenous construction with
zero Nat). Full \texttt{lake build}: 3631 jobs, no \texttt{sorry}.

\paragraph{Contributions.}
\begin{enumerate}
\item The complete empirical trail E1--E18, each experiment documented
individually (\S\S3.1--3.18): design reason, procedure, result, and
significance, as a reproducible artifact (20 scripts, 15 result files,
one-command reproduction), with early first-successes recovered from the
session log (2026-07-29/07-30 prefigurations; perfect generalization
2026-08-15 04:01; reproduction 2026-08-16 14:54, per-seed identical).
\item The Riemann direction (C011--C025, 2026-08-06 visualization $\to$
2026-08-12 formalization and publication) and the homework-exercise series
(2026-08-13 $\to$ 08-15, 25 reports, each with Lean 0-\texttt{sorry}).
\item The five structural laws extracted from the observations, and their
mapping to formalized framework claims (R136/R138/R143/R150/R151/R163).
\item The logical completeness theorem (decoupling operator $D$ = R1
exclusion + R2 cancellation + R3 layering $\Longrightarrow$ completeness at
any order/element/subject-object, verified 2026-08-15 09:14--09:47 with the
two-symmetry-groups corollary I7v) and the three mapping-judgment theorems
(symbol-norm / presentation-legality / intuition-precision, formalized in
\texttt{MappingJudgmentTheorems.lean}, 0 \texttt{sorry}, 2026-08-16).
\item The Lean formalization (150+ Toolkit files, claims R001--R232, 0
\texttt{sorry}, 3631 jobs).
\item A timestamped evidence chain (session records, 2303-turn task log, git
history, file timestamps) and a methodology for recording research trails
under provenance pressure.
\end{enumerate}

\section{The starting observation: the 0/1 OOD dichotomy}

\subsection{Timeline anchors}

\begin{itemize}
\item \textbf{2026-07-29} (session start): first \emph{imperfect}
generalization. A pure next-token transformer memorizes 2-digit addition
100\% (3-digit 0\%). At 22:39:12 (UTC+8) the relation-leave-one-out
generalization is confirmed (training 5 relations, leaving 1: $\le$
generalization 81\% $\to$ 90\% with a new architecture). This is the trail's
first generalization observation. It is \emph{imperfect} --- not all
operations generalize --- but it is already directional: it shows that
generalization is possible without the ``necessary number of steps'', and it
determines the exploration direction (which representations generalize, which
collapse). The same evening produced the first per-token analysis (22:41,
revealing the model's ``reverse'' strategy --- the origin of the token-level
reporting discipline) and the four-question-type system (judgment /
choice / fill-in / computation, 22:49).
\item \textbf{2026-07-30}: early first-successes that prefigure the E-series.
  Cross-class generalization (arithmetic/relations $\to$ logic) 90\% same-class
  / 42.3\% cross-class (00:03); logic-operator leave-one-out (leave $\equiv$
  90\%, leave $\to$ 92\%, leave $\leftrightarrow$ 82\%, 00:28); base-change
  equivalence $\to$ addition generalization 81\% (00:42); the judgment
  paradigm (soft 84\% / vq 71\%) and the parallel-token verification (leave
  $\le$ breaking the counterpart axis: 55\% vs 90\% with complete counterparts,
  05:30); the $\le$ generalization 47\% $\to$ 79\% when logic-gate co-training
  completed the parallel counterparts (08:45). These are the first
  instantiations of what E-series re-runs later formalized as E1/E3/E7--E13.
\item \textbf{2026-08-07}: the basepoint-movement theory is proposed by the
user (thinking document \emph{On the Riemann Conjecture}: the lambda
origin drifts on the circle of radius 1 around 1; iteration continues
along the drift path; the complex plane's origin is suspect --- the two
axes' origins are by definition very likely completely different). The
\texttt{relative-recursion} project starts the same morning (10:30) with
the basepoint question as its working problem, and its first Lean
no-\texttt{sorry} (C001 semiconjugacy) lands at 11:31.
\item \textbf{2026-08-09 01:27}: the user's base(base,base) concept
(ChatGPT conversation ``Lambda calculus and Church numerals'') --- the
lambda basepoint self-application presupposes an operation concept
distinct from the basepoint ($\mathrm{lambda}(\mathrm{base},
\mathrm{base})$), while the intended concept has none
($\mathrm{base}(\mathrm{base},\mathrm{base})$). Same day it becomes the
\texttt{pure-relative} branch of the Lean formalization (claim P001
PureBase; \texttt{C0NoGo}, \texttt{SelfRef}, \texttt{MinimalRuleAudit}).
\item \textbf{2026-08-15 04:01}: first \emph{perfect} generalization: all 19
operations learned, OOD 0.000 $\to$ 1.000 per epoch (reaching 1.000 at epoch
14). The E-series re-runs in the v5 token system follow at 06:48--09:00
(experiment day E1--E13).
\item \textbf{2026-08-16 14:54}: the user asked to reproduce the transformer
experiments; E12 was reproduced 15 training runs per-seed identical, plus a
third in-package run (self-contained package, 08-16 15:04). E14--E18
followed the same day.
\end{itemize}

\subsection{The daily record, 2026-08-01 $\to$ 08-07}

Every day in this window carried work --- trainer design, token-system
theory, tooling, and the launch of the basepoint-movement research. The
window has one unifying theme, stated by the user on 08-02 and verified
by every subsequent day: \emph{can the model's trained intuition be
induced, rather than corrected?} The user's principle: ``what you should
do is make the default methods that are trained into the model succeed''
--- the infrastructure must route the model's native intuitions, not
replace them with new behavior. The R-series prompt experiments (08-03)
are the direct test: pseudo-code headers, code anchors, and few-shot
templates each \emph{induce} better execution (R8 97 vs R11 30 without
them), while negative prohibitions and over-constraint are shown to
\emph{mislead} the intuition (prohibitions suggest the method; tight
constraints collapse). The 08-07 trigger-rate work (0\% $\to$ 40.9\%)
tests the precondition of induction --- the skill must fire before it can
steer; and C001's reuse of mathlib's \texttt{Semiconj.iterate\_right}
instead of re-deriving is the induction of the model's retrieval
intuition. The same theme continues after this window: the CoT-before
tool hint (lean-cot-assist, 08-07/08-09) that lowers the 27B's algebraic
expansion load, and the R-series re-runs on the local 27B after 08-14.
The record is explicit about what was verified: induction works (prompt
shape measurably steers execution), induction has limits (negative
instructions invert), and induction is toolable (trigger description,
hijack-to-correct, hint injection).

The day-by-day record (from the session log, 08-01: 11 entries, 08-02: 31,
08-03: 17, 08-04: 8, 08-05: 13, 08-06: 5, 08-07: 21): The
day-by-day record (from the session log, 08-01: 11 entries, 08-02: 31,
08-03: 17, 08-04: 8, 08-05: 13, 08-06: 5, 08-07: 21):

\begin{description}
\item[08-01 --- trainer deepening and token completeness (11 entries).]
Generalization supervision (train/val split) landed; the user corrected
the input format: ``you are wrong --- the input must be an element
sequence, where every element is composed of its own logic tokens'' (the
raw-string input was rejected). A triple-search harness (three parallel
agents over module-A $\times$ module-B $\times$ router) found the best
combination (layer=tensor, transit=overlap, router=scalar\_only); full
triple enumeration gave 39.8\% overall generalization. The strict $a+b$
extrapolation scored 0/52 = 0\% --- the user's diagnosis: ``parallel
element / parallel token training is insufficient; from the start we must
train addition/subtraction/multiplication, bit operations, base changes,
powers, all logic operations, all question types, with the best method
together''. The user's three completeness questions fixed the token
design: (i) is there a question type for every token's generalization? (ii)
should value not be a scalar? (iii) is every derive traced to baseloop
with training samples designed? --- resulting in 25 zero-coverage token
question types, value dtype scalar $\to$ int/classification, and
definition propositions 16 $\to$ 46 (42/42 complete). The meta-layer was
acknowledged: let the model pour information into itself during training.
\item[08-02 --- tooling infrastructure and the intuition-adaptation
beginning (31 entries).] Llama-server management skill; diagnostic-tool
plugin system (12 plugins); the technology choice ``.sh only for light
orchestration, resident services in JS'' (sh-infrastructure skill). The
user's central correction launched the trail-long intuition-adaptation
methodology: ``what you should do is make the default methods that are
trained into qwythos succeed'' (instead of correcting the model's
behavior) --- the principle that infrastructure must route the model's
native intuitions, not replace them. Prefix-cache-hit analysis and
statusline data integration into claude local; meta-logic/meta-method/
meta-rule research (complete meta-concept landscape, 8 files); the
architecture workflow ``research $\to$ design $\to$ implement'' with
fact-check and judgment-check at every step.
\item[08-03 --- skill writing method and prompt iteration (17 entries).]
The user's corrections became rules: pseudo-code must use code commands
(``run fn'', not ``execute''); negative prohibitions act as hints
(``prohibitions directly violate our experience; they don't prevent, they
suggest the method'') --- rules must be positively phrased. The R-series
prompt experiments (R8: 97 requirement lines with pseudo-code header +
code anchors; ablation R9=46 / R10=63 / R11=30; R12=57 / R13=72 with
few-shot templates) established the prompt anatomy: natural-language task
name first (token direction), then pseudo-code activation, few-shot
templates as format anchors, no over-constraint. Self-generated prompt
research: have the 27B iterate a bad prompt into R13 quality (the
evaluation script was rejected as answer leakage). Deterministic errors
are treated as 100\% errors: hijack the wrong command to the correct one,
or instruct the model which alternative method to use.
\item[08-04 --- technology-stack switch and architecture workflow
(8 entries).] opencode installation failure diagnosis (CachyOS mirror
404s); the switch to opencode Go via a ccproxy (Rust bridge cut from
cc-switch, anthropic-to-openai mode, dual providers); statusline migrated
verbatim from \texttt{.claude} (``I didn't ask you to change the color
scheme or the logic''); the dsv4-flash 1M context displayed as 200k
fixed at the gateway (n\_ctx 1000000); skill-activation research led to
skill-matcher.js (automatic text matching, no AI model in the loop);
G-layer/SKI error summary document.
\item[08-05 --- skill desensitization and trigger-rate optimization
(13 entries).] Switch to opencode dsv4 with the new ccproxy (Rust bridge
with first-byte timeout failover); token-system skill improved; steering-8b
local-run exploration (causal diffusion LM --- not supported by
llama.cpp/GGUF; transformers + bitsandbytes 4-bit path prepared); the
harness results published to GitHub (harness-for-all) after the user
deleted the first 37-file release and required only two skills, fully
desensitized (no claude traces, no dynamic-workflow traces); all 7 skills
desensitized; RTX PRO 5000 evaluation for DeepSeek-V4-Flash-MTP-DSpark-MLX
(conclusion: cannot run --- MLX is Apple-only and the model is a 304B
MoE).
\item[08-06 --- bridge startup and the Riemann visualization (5 entries).]
bridge/start.sh + opencode.sh optimization (stale-socket cleanup, flock
inheritance fix); the ``170hx unlocks 80G VRAM'' claim verified (true
mechanism: CMP 170HX is a full GA100 die; stable community range 10G
$\to$ 40G); and the Riemann zeta / prime-position visualization project
(17:10 request $\to$ 17:55 delivery): 52/52 numerical verifications,
Lehmer $\pi(x)$ to $10^{13}$, browser-tested interactions, complex-base
(Gaussian-integer) support with 14 valid bases (2800 round-trip tests
pass), and persisted prime/zero data (664{,}579 primes; 138{,}067 zeros;
Riemann--Siegel correction 0.1 $\to$ $9\times10^{-4}$).
\item[08-07 --- basepoint-movement research launch + first Lean
no-\texttt{sorry} (21 entries).] The \texttt{relative-recursion} project
started (10:30) with the basepoint question as its working problem (the
user's thinking document \emph{On the Riemann Conjecture}: basepoint
drift, iteration along the drift path, the suspect origin); the C001
semiconjugacy formalization completed with no \texttt{sorry} at 11:31
(first Lean commit 10:53, reusing mathlib's
\texttt{Semiconj.iterate\_right}); the toolchain was built the same day
(allowlist plugin system with hot reload and masking; websearch with
--author and on-disk results; the natural-number-origin research
discipline written into the repository rules; T1--T3 verification tools;
info\_units four-layer information elements; the discipline-guard hook;
the Mathematical Research Protocol skill; description trigger-rate
diagnosis 0\% $\to$ 40.9\%).
\end{description}

\subsection{The two-pole pattern (E12)}

\begin{center}
\small
\begin{tabular}{llll}
\toprule
Group & Training samples & OOD target & OOD (45 runs) \\
\midrule
1a baseline & succ$_a$ 250 & succ$_a$ & \textbf{32/32} $\times$ 9 \\
1b zero-shot & succ$_a$ 250 & \textbf{succ$_b$} & \textbf{0/32} $\times$ 9 \\
2 bridging & succ$_a$ + equivalence 270 & \textbf{succ$_b$} & \textbf{32/32} $\times$ 9 \\
2$'$ control & succ$_a$ + equivalence 270 & succ$_a$ & 32/32 $\times$ 9 \\
4 symmetry & succ$_a$/succ$_b$ mixed 252 & mixed & 32/32 $\times$ 9 \\
\bottomrule
\end{tabular}
\end{center}

All 45 runs reach training accuracy 1.000; the OOD two-pole (0/32 vs 32/32)
never deviates. Token-level analysis locates the collapse: the 0-pole fails at
the notation-branching position (input head passes 32/32; the failure is the
notation fork), a mechanism fingerprint invisible to aggregate accuracy.

Interpretation: training intuition is notation-bound statistics; a 1.000
training accuracy does not transfer to an untrained notation; OOD 1.0 only
shows that \emph{statistics extrapolate within that representation}. The
equivalence-declaration samples (pairwise declarations, anchoring errors
irrelevant) bridge the two notations --- declarations make the information
point at the correct structure.

\section{The Riemann direction and the basepoint-movement theory (2026-08-06 $\to$ 08-12, before the experiments)}

Before the experiment series, the trail ran a parallel formalization branch:
the Riemann direction (C011--C025). Its intellectual spine is the
basepoint-movement theory, proposed by the user on 2026-08-07 (the thinking
document and the \texttt{relative-recursion} project session), formalized
in Lean in the same project the same day (C001, 11:31), carried into the
Riemann direction on 2026-08-12 (basepoint drift under projection,
C011), and finally witnessed empirically by the training experiments
(E11 anchor binding, 2026-08-15 07:43). It starts from the visualization project
(2026-08-06 17:10 request $\to$ 17:55 delivery, within 45 minutes: the
Riemann zeta landscape, the critical line $\sigma=1/2$ highlighted with
nontrivial zeros plotted, prime positions $\le 10^8$, $\pi(x)$ approximations
$\le 10^{13}$, bases 2--36 and Gaussian-integer complex bases; the complex-base
extension and the critical-line-only layout followed at 17:59--18:09, and the
persisted prime/zero data at 18:35--18:46, 138{,}067 zeros with a
Riemann--Siegel correction fix bringing the zero-location error from 0.1 down
to $9\times10^{-4}$). The user's thinking document \emph{On the Riemann
Conjecture} (written around 2026-08-07, between the visualization and the
formalization) records the intuition behind the direction: primes have no
construction method (all other numbers have many); irrationals are the
remainder of a higher-dimensional structure after projection lost part of its
carrying structure; and the pair $1$ vs $1/a+1/b$ (translation vs inversion)
leads to the \textbf{basepoint-drift conjecture --- the origin of the
basepoint-movement theory, 2026-08-07}: the lambda-calculus origin is not
fixed; it moves on the circle of radius 1 around 1, and when the drifted
basepoint's successor iteration is projected back to our number line, the
projection lands on primes. The document adds a second hypothesis: if the
basepoint drifts, iteration continues \emph{along the drift path} --- then
primes land exactly on the translated integer positions, and one can
observe the integer points on that axis. The basepoint-movement research
itself had started the same day in the \texttt{relative-recursion} project
(session 10:30; the project's working question: structures without a
privileged zero, choosing different basepoints interpreted as ``0'') with
its first Lean no-\texttt{sorry} at 11:31 (C001 semiconjugacy); the
thinking document refers to this as ``the basepoint-drift experiment Lean
formalization I did earlier''.

On 2026-08-09 01:27:22 (UTC+8) the user's earliest basepoint-formalization
concept landed in a shared ChatGPT conversation (``Lambda calculus and
Church numerals'', archived in the artifact): the lambda-calculus
basepoint self-application iteration \emph{presupposes a concept distinct
from the basepoint} --- the self-application operation itself, written
$\mathrm{lambda}(\mathrm{base},\mathrm{base})$; but in the intended
experimental concept there is no such operation, written
$\mathrm{base}(\mathrm{base},\mathrm{base})$ --- the same primitive
occupies the operator, operand, and basepoint positions. This
$\mathrm{base}(\mathrm{base},\mathrm{base})$ design became the
\texttt{pure-relative} branch of the formalization the same day (claim
P001, \texttt{PureBase}; Lean files \texttt{C0NoGo}, \texttt{SelfRef},
\texttt{MinimalRuleAudit}), the earliest Lean embodiment of the
basepoint-movement idea: the basepoint is not a distinguished object but a
role. The suspicion about the complex plane's
origin belongs to the same period and the same direction: the two axes'
origins, by their definitions, are very likely \emph{completely different}
--- the origin is not trustworthy. In the 08-12 session the user makes this
explicit and, from then on in the session log, refers to the complex
plane's real axis as the ``fake real axis'' (\emph{jia shi shu zhou}).

The Lean formalization session ran 2026-08-12 00:00--02:54:
\begin{itemize}
\item 00:00--00:12: six prime-drift position observations formalized
  (\texttt{PrimeDriftPositions.lean}, no \texttt{sorry});
\item 00:16--00:21: the user's projection hypothesis formalized as the true
  complex axis --- a two-dimensional rotation algebra $a+bJ$, $J^2=-1$, where
  $\sqrt{-1}$ exists \emph{before} projection (\texttt{ComplexAxis.lean});
\item 00:23--00:27: the user re-states the 08-07 basepoint-drift
  conjecture against the freshly built complex axis and sharpens it: since
  the complex plane's axis carries the naturals, the lambda basepoint has
  moved onto the complex axis --- \emph{the basepoint has drifted, its
  position is not the origin but some transform of $i$} (user, 00:23,
  quoted in the session log). Then (00:25): ``that so-called real axis is
  very suspicious --- I suspect it is not a real number axis at all, but
  another fake complex axis''. This is where the ``fake real axis'' term
  enters the log; the user uses it consistently from then on. Both
  statements were formalized the same night: the basepoint drift (the true
  position is unobservable under projection) and the fake-real-axis theorem
  (the real axis is a sub-algebra $\{\langle a,0\rangle\}$ of the same
  structure, its position unobservable --- the ``origin illusion'');
\item 00:28--00:29: C011 + C012 filed as the first claims of the direction;
\item 00:53--02:04: the chain C013--C025 (projection frame, prime circles
  C012--C018, critical-line geometry C019--C022, prime-circle products
  C023--C024, Euler-product convergence with the zero-free region C025);
\item 02:54: session archived with C011--C022.
\end{itemize}

\paragraph{The claim chain C011--C025, item by item.} The direction is a
single intuition chain, each link formalized in Lean (all PROVED, novelty
KNOWN --- restatements of classical mathematics; the record states
explicitly, per the user's instruction, that no claim of proving the
Riemann hypothesis is made and mathlib's official statement remains
unproved):

\begin{enumerate}
\item \textbf{C011 (projection construction of the complex plane)}: the
complex axis is the projection of a two-dimensional rotation algebra
$ComplexAxis = \{\langle a,b\rangle\}$, $J=\langle0,1\rangle$, $J^2=-1$;
$\sqrt{-1}$ exists before projection; basepoint drift under projection
(the basepoint's true position is unobservable --- the user's ``fake real
axis'' is the theorem: the real axis is a sub-algebra of the complex
rotation algebra, not an independent axis; its position is unobservable
under projection --- the ``origin
illusion''). Filed 08-12 00:29 (commit 133f353).
\item \textbf{C012 (prime drift positions)}: primes land on translated
integer points of the drift axis; six positional observations
(integer decomposition $n\cdot p$ prime $\Rightarrow n=1$, irrationals
have no integer multiple, harmonic identities) formalized in
\texttt{PrimeDriftPositions.lean}, 00:00--00:12, commit 7fdcb91.
\item \textbf{C013 (Riemann direction in the projection frame)}: inversion
drift, partial sums, and the critical line as 15 theorems in the
projection frame, 00:53--01:02 (commit 7c0feb8).
\item \textbf{C014 (sums of two squares, C015 curling)}: a prime
$p \not\equiv 3 \pmod 4$ is the norm $a^2+b^2=p$ of some complex point
(citing mathlib's Fermat two-squares theorem); the inversion
$\mathrm{recip}\, z = \overline z / |z|^2$ curls infinity back to
finiteness (for every $\varepsilon$ there is $R$ with $|z|>R \Rightarrow
|\mathrm{recip}\, z|<\varepsilon$), is an involution, and satisfies
$|\mathrm{recip}\, z| = 1/|z|$ --- the geometric mechanism of analytic
continuation. C014 commit 20df6b6 (01:13--01:21), C015 commit de3e966
(01:38).
\item \textbf{C016 (lattice points on a circle)}: the orbit structure of
integer points on prime circles, uniqueness via the Gaussian-integer UFD
($\mathbb Z[i]$ is a Euclidean domain), 01:42 (commit e235d09).
\item \textbf{C017 (prime circles and conjugate pairing)}: a prime
$p\equiv1 \pmod 4$ circle carries exactly one orbit of 8 lattice points
(up to units $\pm1,\pm i$), splitting into conjugate Gaussian primes;
01:55 (commit edbc099).
\item \textbf{C018--C022 (critical-line geometry)}: the critical line as
the inversion--translation duality center (C019), the critical-line circle
through 0 and 2 (C020), the reflection-condition equivalence
(line $\iff$ reflection condition $\iff$ circle), and the functional-
equation symmetry restated algebraically --- explicitly not a claim about
zero positions.
\item \textbf{C023--C024 (prime-circle products and geometric splitting)}:
products of prime circles and their decomposition into conjugate pairs
under rotation.
\item \textbf{C025 (Euler-product convergence and zero-free region)}: for
$\mathrm{Re}(s)>1$, $\prod_p (1-p^{-s})^{-1} = \sum_n 1/n^s$ = mathlib's
official $\zeta(s)$ (via
\texttt{eulerProduct\_completelyMultiplicative\_tprod} +
\texttt{Complex.summable\_one\_div\_nat\_cpow}); $\mathrm{Re}\ge1$ is
zero-free.
\end{enumerate}

The honest boundary is recorded in the claim chain itself: every claim is
labeled KNOWN (restatement), the Riemann hypothesis and the twin-prime
conjecture are neither proved nor touched, and the positional claim on
zeros (all nontrivial zeros on $\mathrm{Re}=1/2$) is not part of any
claim. The methodology record (token economy: $\approx$700k tokens, 1009
requests, 99.2\% context-cache retransmission) is reported as an
efficiency case study, with the single data point that after a sleep
interval the intuition compacts --- recorded, not concluded.

\paragraph{First-completion times (Riemann direction).}
The visualization was delivered 2026-08-06 17:55 (request 17:10, UTC+8), with
the complex-base extension completed 18:09 and the persisted data 18:46. The
first Lean no-\texttt{sorry} of the whole trail is C001 (semiconjugacy),
completed 2026-08-07 11:31 (first Lean commit 10:53). The Riemann-direction
chain C011--C025 was completed 2026-08-12 00:00--02:54; the technical record
and the full research paper were dated and published 2026-08-12 (technical
record DOI 10.5281/zenodo.21896990, full paper 10.5281/zenodo.21897167, and
the earlier DOI 10.5281/zenodo.21896345 added to the main paper at the
user's instruction 09:43--09:44, after the user corrected the model's
referent choice at 09:44: ``no, the Riemann formalization's main paper'').
The basepoint-drift result inside the Riemann direction is the
first-completed instance the user cites for the ``first completion counts''
rule. The user's instruction at 2026-08-13 03:05 --- ``revise the Riemann
paper wording: every claim about 'proving' is what DeepSeek insisted on'' ---
corrected the record's attribution: the unproved statements are the API
model's insistence, not the user's claim.

\section{The homework-exercise series (2026-08-13 $\to$ 08-15)}

Parallel to the experiments, the trail produced a series of ``homework
exercises'' for the framework's foundation paper: re-observations of classical
problems through the Pat lens, each with its own Lean formalization anchored
only to framework theorems. The series ran 2026-08-13 14:19 through
2026-08-15 (25 reports: exercises I--XXII + Pat0 overview + summary).

\textbf{Exercise I (P vs NP)} started the series at 2026-08-13 14:19 (session
``PAT-view P vs NP paper and Lean verification''). The user's corrections
steered it through three iterations in 27 minutes:
\begin{itemize}
\item 14:24 --- positioning correction: ``not the computation chapter of the
  psyche paper, it is a homework exercise for the foundation chapter''; the
  exercise expands claims R152/R153 (P vs NP in the Pat frame);
\item 14:31 --- scope correction: ``do not over-enumerate, meaningless; under
  this framework we trade speed for fuzziness and recover precision with
  structure''; the initial 9-theorem enumeration was deleted, leaving the 4
  mechanism theorems that carry the single thesis: NP's Pat semantics = ``use
  fuzziness for speed (existence via phase-locking extrapolation R150), use
  structure to recover precision (verification table: witness $\iff$ table
  entry)''; the full build stayed green;
\item 14:45--14:46 --- the user asked for a detailed cross-reference document,
  and that all P vs NP material live in the exercise's own subfolder
  (a self-contained package pattern that all later exercises follow);
\item 14:48 --- citation instruction: ``cite as many living mathematicians,
  philosophers, linguists and scientists as possible'', with the discipline
  that every citation is verified via web search (no citations from memory).
\end{itemize}

\textbf{Exercises II--IX} (each: paper ZH/EN + Lean + claim + README) were
written during 2026-08-13 in parallel sessions; at 18:26 the user asked to
publish all exercise papers and Lean files. Ten DOIs were published
(2026-08-13 18:26--18:37, 43 files, all HTTP 201). At 19:01 the user
extended the requirement: ``up to exercise 14, all must have double-blind
publication structure''; exercises X--XIV had just been created by the
parallel sessions, and the double-blind packages (15 at that point) were
built and published within minutes.

\textbf{The exercise list (topic / anchor / claim).} Each exercise is
self-contained (paper ZH+EN, Lean snapshot, claim YAML, README) and
anchors only to framework theorems:

\begin{itemize}
\item \textbf{I: P vs NP in the Pat view (R157, 08-13 14:19--14:46)} ---
expands R152/R153: NP's Pat semantics = ``use fuzziness for speed
(existence via phase-locking extrapolation R150), use structure to recover
precision (verification table: witness $\iff$ table entry)''; 4 mechanism
theorems, 0 sorry; the cross-reference document added 14:45--14:46.
\item \textbf{II: the 1 of each number domain (R158)} --- $\sqrt2/2$,
complex, arbitrary $n$-ary, transcendental numbers: each domain's ``1''
as its basepoint; 8 theorems, 0 sorry (PatNumberOnes.lean, anchored to
R154's $\sqrt2/2$ projection).
\item \textbf{III: Riemann conjecture and twin primes revisited (R159)}.
\item \textbf{IV: trigger-stripping experiment (self-reference collapse
initiation)} --- the mechanism by which unpaired self-reference collapses.
\item \textbf{V: four interlock as the minimal self-consistent interlock
(R160)}.
\item \textbf{VI: interlock growth step and projection detachment
(R161/R162/R163)}.
\item \textbf{VII: five forces domain-by-domain Pat observation (R164)}.
\item \textbf{VIII: Pat-native group representation theory (R165)}.
\item \textbf{IX: physical space interlock structure (R166)}.
\item \textbf{X--XIX: the seven millennium/classical problems re-observed}
--- Navier--Stokes (R167), Hodge (R168), BSD (R169), Goldbach (R170),
parity and primes (R171), twin-prime conjecture (R172), twin-prime gaps as
an oscillation-period axis (R173), Collatz (R176), ABC (R179), odd
perfect numbers (R180).
\item \textbf{XVIII: the real axis in the Pat mapping architecture (R178)}
--- the real axis as a projection of the complex rotation algebra.
\item \textbf{XXI: continuum = closure of the Pat grid} and
\textbf{XXII: composite numbers = multiple projection superpositions of
primes} --- the two summary observations connecting the series back to the
continuum construction (R150/R151/R163).
\end{itemize}

Every exercise carries the structure ``exercise / topic / Lean / DOI''
(per the user's 08-14 12:47 re-qualification) and the declaration that it
represents no mathematical conclusion, only an observation report based on
Lean formalization; every Lean file compiles with 0 \texttt{sorry}.

\textbf{Exercises XV--XXII} continued on 2026-08-14 (R167--R180: Navier--
Stokes, Hodge, BSD, Goldbach, parity and primes, twin primes, Collatz, the
real-axis mapping architecture, ABC, odd perfect numbers). On 2026-08-14
12:47 the user re-qualified the whole series: every exercise must carry the
structure ``exercise / topic / Lean / DOI'' and declare that it
\emph{represents no mathematical conclusion, only an observation report
based on Lean formalization}. The series converged on 2026-08-15 with 25
reports, all with Lean 0-\texttt{sorry} formalizations.

\section{The user's theoretical findings: the token-system theory (2026-07-29 $\to$ 08-14)}

Before and during the E-series, the user built a theoretical layer about
what training intuition is and how representations should be organized.
The E-series is its empirical testing ground; the framework formalizes its
surviving parts. This section records the findings themselves (verbatim,
with timestamps), each with its later landing point in the experiments or
the formalization.

\begin{enumerate}
\item \textbf{Induction and deduction are construction and intuition}
(2026-07-29, thinking document): ``I understand that feeling and expression
are like induction and deduction --- a pair of inverse structures;
induction relies on construction, deduction relies on intuition. This
means that even within the same expression system, receiving and
expressing, feeling and executing, are very likely \emph{not to use the
same token} even in the same modality.'' Landing: the symbol/concept layer
separation of the token system and E2 (symbol mapping direction).
\item \textbf{The problem is also an element} (2026-07-29 22:46): ``the
problem/operation and the data are isomorphic --- both are elements, both
are logic-token vectors.'' Landing: the four-question-type system
(judgment / choice / fill-in / computation), where the question itself is
an element sequence; later the whole sample-synthesis architecture.
\item \textbf{Formal reasoning is the position of tokens, not tokens}
(2026-07-29 18:38): ``formal reasoning is essentially not tokens, but
\emph{positions} of tokens; induction is the discovery of positions,
deduction and abduction are the use of positions; and a position can be a
coordinate in a matrix, not necessarily a graph relation'' (the model had
wrongly narrowed positions to graph relations; the user corrected:
``who said position must be a graph relation? Can't it be a position in
some matrix?''). Landing: the position-value representation study (E6) and
the chain-length representation.
\item \textbf{The knowledge starfield is a high-dimensional Hilbert space,
not a tree} (2026-07-30 correction): ``my understanding of logical
abstraction factorization was completely wrong; the real knowledge
starfield is a high-dimensional Hilbert space, not a tree''; and ``$=$/$\equiv$
is a set of multiple logic tokens in the space (the intension and
extension of a concept); independent points are the properties (logic
tokens) that compose things''. Landing: the logic-token design discipline
--- a concept is a token set, an element is a composition, and a
``collision'' (two concepts sharing a token pattern) is an incomplete axis
set, not a bug.
\item \textbf{Token cancellation} (2026-07-30): ``only tokens that all
input elements share \emph{with the same value} are cancelled'' (the
user rejected the model's ``cancel 0'' shortcut: ``why cancel 0? My
requirement was: if all elements at the input layer share the same token,
cancel that token; I never said cancel 0''). Cancellation is set
intersection --- pure compilation, no dimensional degeneration in high
dimensions. Landing: the cancellation/collapse mechanics of the v2--v5
training pipelines.
\item \textbf{The sample-size information bound (Mastermind steps)}
(2026-07-29 22:20--22:24): ``data synthesis should guarantee that every
logic token pairing is trained enough --- like the four-color vs
eight-color guessing game: you cannot limit both to the same number of
steps''; ``the minimum guessing-game steps: given eight colors and four
positions, each guess reports how many colors and how many positions are
correct'' --- the sample count has an information-theoretic floor.
Landing: information-gain sampling (the halving experiment, 22:35; the
first generalization 22:39:12) and later the E-series sample budgets.
\item \textbf{Weight compression} (2026-07-30 08:17): ``I understand this
is essentially training more fully with fewer logic tokens --- what we are
doing is equivalent to compressing weights'' (registering the logic gates
AND/OR/NOT/XOR as logic tokens: the gates share the boolean polarity
axes, so one token set covers eight operators). Landing: the 47\% $\to$
79\% $\le$-generalization result when logic-gate co-training completed the
parallel counterparts (07-30 08:45); the ``compression via shared axes''
principle behind the token vocabulary.
\item \textbf{Parallel logic tokens, not parallel elements}
(2026-07-30 05:30): ``once all logic tokens of an element are
sufficiently trained through parallel logic tokens, the element
generalizes even if never trained''. Landing: the 55\% vs 90\% two-pole
(leaving $\le$ with a broken counterpart axis vs complete counterparts);
the E-series two-pole designs.
\item \textbf{The convergence phenomenon as an overfitting question}
(2026-08-14): the observed token-training convergence (loss 6.59 $\to$
1.53) must be examined for whether it is overfitting; the author's
exploration starts from ``what changes under pure-intuition overfitting''.
Landing: the overfitting discussion in the observation report and the
E-series recovery judgments.
\item \textbf{Directionality as a primitive notion} (2026-08-01):
direction must have a precise mathematical/logical definition
(directed/undirected), because directions determine which pairs can be
declared (R136: directions must be declared in pairs $\{d,-d\}$, twice
declared = asymmetry). Landing: the paired-declaration law and E18
(paired 5/9 vs one-way 1/9).
\item \textbf{Symmetry as dimensionality reduction, the restoration
point = midpoint} (2026-08-15 06:00): ``symmetry is dimensionality
reduction''; the restoration point (midpoint) of a symmetric pair is a
generation rule. Landing: R163 (endogenous continuum construction) and
E14 (symmetry center not at 0; center drift extrapolates).
\item \textbf{Orthogonality is itself a symmetry} (2026-08-15 09:14) and
\textbf{two symmetry groups localize a direction} (09:19, I7v): ``Pat0's
declaration must be two groups of symmetries; any thought experiment based
on a single dimension is wrong''. Landing: the 09:16--09:19 verification
(double declaration 5/5 stable; direction localization 16/16 $\times$ 3
seeds) and the logical-completeness project (Section~8.1).
\item \textbf{The logical completeness theorem} (R031, C1--C7): the
decoupling operator $D$ (R1 exclusion + R2 cancellation + R3 layering)
$\Longrightarrow$ completeness at any order/element/subject-object:
perfect OOD. Landing: Section~8.1 (verification project 2026-08-15
09:14--09:47).
\item \textbf{Mechanics symmetries come in interlocked orthogonal pairs}
(2026-08-16, R233, 29 theorems): each degree of freedom is a four-phase
interlock pair; the symplectic structure $J^2=-I$; the full embedding is
12 pairs $= \mathbb C^{12}$; the gravitational direction is asymmetric
(structure layer has no privilege, action layer is one-way). Landing: the
R233 formalization (29 theorems, 0 \texttt{sorry}) and the E18 mechanics-
token experiment (paired declaration 5/9 vs one-way 1/9).
\item \textbf{The five-layer token system B/C/S/G/P} (2026-07-29 $\to$
08-16): basepostulate (B) / concept (C) / symbol (S) / grammar (G) /
presentation (P); the layers have distinct roles --- B/C are the meaning
layer (what things are), S/G are the expression layer (how things are
written and arranged), P is the concrete notation (symbols with
precedence and associativity). A token is one thing; the JSON files are
projections of it. Landing: the whole v5 tokenizer and the three mapping-
judgment theorems below.
\item \textbf{The three mapping-judgment theorems} (2026-08-16
16:17--16:23, the user's correction re-located them): how to \emph{judge}
the multi-to-multi mappings between layers --- (i) \textbf{symbol-norm
theorem}: the judgment of the multi-to-multi mapping from symbols to logic
and intuition (every target must be resolvable; ambiguity converges by
context; unresolvable = non-norm); (ii) \textbf{presentation-legality
theorem}: the judgment of the mapping from logic and intuition to
presentation (every presentation target must be resolvable; numeric
coincidence without declaration is illegal); (iii) \textbf{intuition-
precision theorem}: the judgment of the mapping from intuition to logic
and presentation (the intuition-token interface must be definitionally
resolvable, must not impersonate a logic concept; precision is judged by
two-pole comparison; polarity-confusion events are the canonical failure).
A fourth theorem (symbol $\to$ presentation) completes the five-layer
coverage. Landing: \texttt{MappingJudgmentTheorems.lean} (5 theorems, 0
\texttt{sorry}, 2997 jobs) with the experiment-coverage matrix
(E6/E4/E5/E12/E13/E15/E14 + R065) --- Section~8.2.
\item \textbf{Symbols are not values; numerals are intuition} (2026-08-15
06:54): ``numeral is itself intuition; the true numeral is the successor
chain from the basepoint; two successors from the basepoint are addition''.
Landing: E6 (structure-bearing representation) and the whole chain-length
representation of the E-series.
\item \textbf{The head paradigm: same-kind/intension-extension/same-value}
(2026-07-31): the head is determined by whether tokens are same-kind,
and within same-kind by the same-value relation of their intension/
extension --- not by hand-crafted heads. Landing: the head-router design
of the v2 trainer and the token-level judgment discipline.
\item \textbf{Number-symbols need both reference and definition}
(2026-08-01): the numeral symbol's signifier vs referent (the symbol
[3] vs the concept 3 differ in token composition); every derived token
needs a definition proposition, not only a derivation. Landing: the
definition-layer discipline (E4/E5: definitions are invisible to
training) and the token-system maintenance rules.
\end{enumerate}

The theoretical layer and the empirical layer are kept separate in this
record: the findings above are the user's claims (each with its verbatim
timestamp), the E-series results are the witnesses, and the Lean theorems
are the mechanical verification. No finding is promoted to a theorem
without a Lean proof; no E-result is promoted to a law without its
two-pole evidence.

\section{The empirical trail: E1--E18, experiment by experiment}

The series is design-driven: each experiment answers the question raised by
the previous one. We document each experiment individually --- design reason,
procedure, result, and significance --- so that the trail can be audited step
by step. Configuration throughout: dim 64, layers 2, epochs 60 (unless noted),
lr $10^{-3}$, seeds $\{0,1,2\}$, pure next-token transformer, token-level
judgment (full-sequence reconstruction; position-level accuracy is not
trusted).

\subsection{E1: extraction-channel ablation (ablate\_channels)}\hfill\textit{(first success: 2026-08-15, experiment day)}

\textbf{Design reason.} Before any law can be claimed from an OOD result, the
judgment itself must be strict: ``training intuition'' is an unobservable
construct, and every claim about it needs a decision procedure. The trail
distinguishes three extraction channels through which a training result can
arise --- the \emph{sample} channel (what the training samples carry), the
\emph{construction} channel (what the representation's construction rules
carry), and the \emph{arrow} channel (what the pre-defined A-layer arrows
carry). The question is which channel is \emph{decisive}: remove it, and the
OOD result must change. This is a completeness-style elimination: a channel
is called decisive only if its removal is witnessed by a two-pole change
(1-OOD $\to$ 0-OOD or the reverse), never by a single observation.

\textbf{Procedure.} Domain: $\mathbb Z/7$ (mod 7), where the multiplicative
inverse exists for every non-zero element --- a non-degenerate point probe
that the additive domain cannot provide. Baseline: all three channels
present. Then each channel was removed one at a time, every combination over
multiple seeds, training accuracy and OOD judged per seed.

\textbf{Result.} Only the sample channel is decisive: removing it changes
the outcome; removing construction or arrows alone does not. The two-pole
evidence: with samples present, OOD succeeds; with samples removed, it
collapses --- regardless of construction or arrow status.

\textbf{Token-level judgment.} All results reported as full-sequence
reconstruction (sample-pair/total + position-level collapse localization);
aggregate accuracy alone is not trusted.

\textbf{Significance and declaration boundary.} Structure still carries
information (ontology: the representation's construction rules and arrows
are what the samples are \emph{about}), but extraction must pass through
samples (epistemology: the training process reads samples, not structures).
The claim is about \emph{which channel decides the training outcome}, not
about structure being useless --- the two statements are separated so that a
reader cannot conflate them.

\subsection{E2: symbol mapping direction (sym\_map\_exp)}\hfill\textit{(first success: 2026-08-15, experiment day)}

\textbf{Design reason.} The token system separates a symbol layer (how things
are written) from a concept layer (what things are). The question: does the
direction of the mapping matter? A bare symbol $\to$ concept mapping crosses
layers; a concept $\to$ concept mapping stays within one layer. If training
intuition is layer-sensitive, the two directions must behave differently;
if not, the layer distinction is presentation-only.

\textbf{Procedure.} Two training setups, same vocabulary, same seed set:
(a) symbol tokens mapped directly to concepts (cross-layer), (b) concepts
mapped to concepts (same-layer). Training accuracy and OOD compared across
seeds.

\textbf{Result.} Symbols must map to \emph{structured} concepts; a bare
symbol mapped directly (without a structured target) is unstable --- the
OOD collapses, while the same symbol mapped through a structured concept
extrapolates.

\textbf{Token-level judgment.} The collapse is located at the symbol-to-
concept junction (the first token after the symbol); the input head passes,
the failure is the mapping itself.

\textbf{Significance and declaration boundary.} This fixes the symbol
discipline used throughout the trail: a symbol is not a value, it points at
a structured concept; ambiguity is resolved by context, never by the symbol
alone. The claim is about mapping \emph{direction and target structure}, not
about the symbol layer being eliminable --- presentation remains necessary
for expression.

\subsection{E3: pure-structure nesting foundation (succ\_nest\_exp)}\hfill\textit{(first success: 2026-08-15 06:48, OOD 36/36; preceded by single-operation failures 06:43 OOD 0/15, corrected to mixed training)}

\textbf{Design reason.} Before any law about anchors, notation, or 0, the
trail must answer the foundation question: is isomorphic nesting --- a chain
whose length is countable and whose parts are shared --- learnable at all?
If nesting is not learnable, nothing that follows can stand. The logic is
layered: (i) numerals as basepoint successor chains, (ii) addition and
multiplication as chain-length operations, (iii) 0 as the empty chain.

\textbf{Procedure.} 4 tokens (zero/succ/addition/multiplication), succ-nested
representation, addition and multiplication with 0 fully enumerated, 5 seeds.
Configuration: dim 64, layers 2, epochs 60, lr $10^{-3}$.

\textbf{Result.} Convergence 1.000 $\times$ 5 seeds; OOD 36/36 (new-range
chain lengths) and 0-combination OOD 60/60 $\times$ 4 seeds. The isomorphic-
nesting foundation is stable.

\textbf{Token-level judgment.} Full-sequence reconstruction passes for all
36/36 and 60/60 samples; the per-position trace shows no collapse at any
chain position (the chain is read left to right and reconstructed
position-for-position).

\textbf{User correction during the run (06:15--06:48).} The first attempt
trained a single operation with small samples and failed (OOD 0/15): ``a
single operation with small samples cannot be learned anyway'' (the earlier
perfect generalization used 4816 mixed samples). The fix --- mixed training
with the target screened --- is itself a recorded correction of the API
execution: the model had to be steered from single-operation isolation to
mixed-representation training before the 36/36 result appeared.

\textbf{Significance and declaration boundary.} Chain-length structure is
learnable and extrapolates within its representation. The claim is scoped
to the succ-nested representation; it says nothing yet about other
representations (that is E6) and nothing about anchors (that is E11).

\subsection{E4: four definitions of 0 (zero\_exp)}\hfill\textit{(first success: 2026-08-15 06:54, OOD 60/60)}

\textbf{Design reason.} 0 can be defined in at least four ways: axiomatic (a
constant), structural (the empty chain), intuitive (the basepoint itself),
or relative (the basepoint of the current frame). The question: does the
\emph{definition} of 0 affect training? If the definition layer is invisible
to training, the four definitions must train equivalently; if visible, they
must diverge. This separates the definition layer from the representation
layer --- a distinction the whole framework depends on.

\textbf{Procedure.} Four definitions of 0 $\times$ multiple seeds, identical
sample structure otherwise (succ-nested numerals, full 0-combination
coverage), dim 64, layers 2, epochs 60.

\textbf{Result.} All four definitions train equivalently: 1.000 training
accuracy, OOD 60/60 across seeds. No divergence in any seed.

\textbf{Token-level judgment.} The 0-combination OOD set (0+$n$, $n$+0,
0$\times n$, $n\times$0, 0$\times$0) passes full-sequence reconstruction
position-for-position; no collapse at the 0 token.

\textbf{Significance and declaration boundary.} Training intuition is
sensitive only to the representation form, not to the definition layer.
The claim is scoped: definitions that \emph{produce the same representation}
are indistinguishable to training; it does not claim all definitions are
equivalent in general --- only that training cannot see the difference.
This is the definition-layer invisibility law (law 5).

\subsection{E5: 0-definition $\times$ succ-design matrix (zero\_succ\_exp)}\hfill\textit{(first success: 2026-08-15, 1.000 $\times$ 24 runs)}

\textbf{Design reason.} E4 varied the definition of 0 alone. The design
matrix completes the picture: 0's four definitions $\times$ succ's two
designs (explicit concept token vs atomic counting) --- 8 combinations. If
both axes are invisible to training, all 8 must train equivalently; any
single divergence localizes exactly which layer is visible.

\textbf{Procedure.} 8 groups $\times$ 3 seeds = 24 training runs, same
sample structure (succ-nested, full enumeration), dim 64, layers 2, epochs
60.

\textbf{Result.} All 24 runs reach training accuracy 1.000; no group
diverges. Succ's explicit/atomic design is as invisible as 0's definition.

\textbf{Token-level judgment.} Verified on full-sequence reconstruction
across all 24 runs (per-run sample pairs reported); no position-level
collapse in any group.

\textbf{Significance and declaration boundary.} The definition layer is
invisible on both axes of the design matrix. Scoped claim: invisibility
holds for the 8 tested combinations under succ-nested representation;
representation changes (E6) are a different axis and are not covered here.

\subsection{E6: structure-bearing vs bare-symbol representation (zero\_symbol\_exp)}\hfill\textit{(first success: 2026-08-15 08:36--08:37; 0-pole: third-representation OOD 0/32)}

\textbf{Design reason.} E4/E5 showed the definition layer is invisible. The
representation layer is the next candidate: is a position-value notation
(implicit place value, nominal position tokens) equivalent to a true numeral
(fully expanded successor chain)? The intuition to test: \emph{structure
must be carried by the representation, not assumed}. A bare-symbol
representation assumes the structure implicitly; a true numeral expands it.

\textbf{Procedure.} Four representations A/B/C/D trained with identical
sample content, then OOD-tested: A = application-chain (true numeral, the
E3 representation), C = chain with value tokens; B/D = position-value /
bare-symbol variants. 30 OOD probes per representation per seed.

\textbf{Result.} A and C: OOD 30/30 (stable). B and D: OOD 0/30 (collapse).
The two-pole is sharp: position-value needs structural support; the true
numeral is a basepoint successor chain and extrapolates.

\textbf{Token-level judgment.} The B/D collapse is located at the position
layer (position 4--14 fails; input head and read-out pair pass) --- a
mechanism fingerprint that distinguishes format-layer failure from
operation failure, invisible to aggregate accuracy.

\textbf{User correction (08:36).} The user's stripping discipline applied
here: ``before a complete ablation comparison, every OOD result must be
assumed to be an intuition path; differences exist only through
comparison'' --- the third-representation 0/32 was only reportable after
the same-representation baseline 32/32 existed (both poles observed).

\textbf{Significance and declaration boundary.} Structure-bearing
representation is necessary for extrapolation; bare symbols collapse. The
claim is about \emph{representation form}, not about the mathematical
content: both representations denote the same numbers, only the
learnability differs.

\subsection{E7: 0-combination holdout (zero\_hole\_exp)}\hfill\textit{(first success: 2026-08-15 07:22, 40/40 + far-OOD 24/24)}

\textbf{Design reason.} A genuine OOD judgment requires that the test set
contain no sample analogous to the training set. Earlier series trained
0-combinations fully (E3--E5); this experiment holds out part of the
0-combinations (some 0$\times n$ and $n\times$0 pairs) and asks whether the
0-absorption law extrapolates to the held-out pairs.

\textbf{Procedure.} Partial 0-combination training (5 of the 0-samples
present), held-out 0-combinations as OOD (40 probes), plus a far-OOD set
(24 probes, 0$\times5\ldots0\times9$ on chains never seen).

\textbf{Result.} Held-out 40/40; far-OOD 24/24. The 0-absorption law
extrapolates from 5 0-samples to 0$\times5\ldots0\times9$.

\textbf{Token-level judgment.} Full-sequence reconstruction passes for all
64 probes; the 0 token is reconstructed correctly at every position it
appears (no position-level collapse).

\textbf{Significance and declaration boundary.} 0 is not a special
exception in the representation; it is chain-length 0, and its rules
extrapolate exactly like any other chain rule. The claim is scoped to the
succ-nested representation; the ``0 is special'' hypothesis (needing
special training) is rejected by the evidence, not by assumption.

\subsection{E8: 0-multiplication vs nonzero multiplication (zero\_vs\_nz\_exp)}\hfill\textit{(first success: 2026-08-15 07:25, 0-multiplication zero-shot 45/45)}

\textbf{Design reason.} E7 showed 0-absorption extrapolates when some
0-samples are present. The sharper question: is 0 special at all? Train
\emph{only} nonzero multiplication (0-multiplication zero-shot), then test
0-multiplication. If 0 is chain-length 0, the multiplication rule must
extrapolate to it without any 0-sample; if 0 needs its own training, it
collapses.

\textbf{Procedure.} Nonzero-multiplication training only; OOD: all
0-multiplication pairs (0$\times n$, $n\times$0, 0$\times$0), 45 probes,
multiple seeds.

\textbf{Result.} 0-multiplication zero-shot OOD 45/45. The chain-length
multiplication rule naturally extrapolates; 0 is chain-length 0, not a
special case.

\textbf{Token-level judgment.} All 45 probes pass full-sequence
reconstruction; the 0 token at both operand positions reconstructs
correctly (no position-level failure).

\textbf{User instruction (06:54).} This experiment exists because of the
user's correction: ``2. Then 0-multiplication must be compared against
nonzero multiplication'' --- the model had to be steered from training
0-combinations inside the sample set to leaving them out entirely.

\textbf{Significance and declaration boundary.} The result is a law-level
claim (0 = chain-length 0) with a zero-shot witness. Scoped: the claim
holds for the succ-nested representation and the multiplication rule as
encoded; it does not generalize to arbitrary operators by itself (E13
tests operation distinction separately).

\subsection{E9: succ shift invariance (succ\_shift\_exp)}\hfill\textit{(first success: 2026-08-15 07:33, succ(1..7) 14/14)}

\textbf{Design reason.} The chain-length reading of numerals assumes that
succ(0)$=$1 and succ(1)$=$2 are \emph{the same} succ applied to different
chain lengths. If that is true, training succ(0)$=$1 alone (2 samples)
must suffice for succ(1..7); if succ were position-dependent, it must fail.
This is the shift-invariance question: is the successor rule translation-
invariant in the representation?

\textbf{Procedure.} Train only succ(0)$=$1 (2 samples) and the 0-identity
samples; OOD: succ(1..7) (14 probes), with a value-zero error-sample control
group. Multiple seeds.

\textbf{Result.} succ(1..7) OOD 14/14. Shift invariance holds: 0$\to$1 and
1$\to$2 are the same succ. The error-sample control shows no difference
(value-zero mistakes do not affect the shift result).

\textbf{Token-level judgment.} 14/14 full-sequence reconstruction; each
succ chain position reconstructs correctly (the succ token at every
position, the numeral at the end).

\textbf{Significance and declaration boundary.} The mechanism of chain-
length extrapolation is shift invariance: the successor rule is
translation-invariant, which is exactly what makes a chain a chain. Scoped:
invariance is witnessed for one succ axis; multi-axis behavior is E10.

\subsection{E10: multi-axis succ independence (axis\_succ\_exp)}\hfill\textit{(first success: 2026-08-15 07:33, same-axis 12/12)}

\textbf{Design reason.} E9 showed one succ axis is shift-invariant. Does
the invariance transfer across axes? Two succ axes (succ$_a$ on axis a,
succ$_b$ on axis b) may share the rule or be independent. The question is
answered by training one axis and OOD-testing the other: if axes share the
succ rule, cross-axis OOD passes; if each axis has its own succ, it
collapses.

\textbf{Procedure.} Dual-axis and single-axis training; single-axis
ablation (train axis a only, OOD axis b). 12 probes per direction,
multiple seeds.

\textbf{Result.} Same-axis 12/12; cross-axis 0/12. Each axis's succ is
independent; a single axis can be ablated away without the other
recovering it.

\textbf{Token-level judgment.} The cross-axis collapse is located at the
first succ$_b$ token (the axis-identity position); same-axis passes
position-for-position.

\textbf{Significance and declaration boundary.} Succ is axis-relative: the
rule is per-axis, not global. This is the first precise statement of
``direction is declared per axis'' --- the empirical seed of the framework's
paired-direction-declaration law (R136). Scoped: independence is witnessed
for two axes; it does not claim a bound on the number of axes.

\subsection{E11: basepoint drift (anchor binding) (basepoint\_drift\_exp)}\hfill\textit{(first success: 2026-08-15 07:43, same-anchor 32/32, drift 0/32; explicit bridging also 0/32)}

\textbf{Design reason.} E9/E10 showed succ is per-axis. The deeper question:
is the \emph{operation} itself bound to its basepoint (anchor)? If
operations are anchor-bound, moving the anchor must break the OOD (drift
0-pole), and adding an explicit bridge must fail to recover it; if
operations are anchor-free, both must pass. This is the reference-frame
question --- the empirical root of the basepoint-relative operator theory.

\textbf{Procedure.} Anchor-0 and dual-anchor training; OOD three ways:
same-anchor (baseline), drift (same operation, different anchor), explicit
bridge (drift + declared equivalence between anchors). 32 probes per
condition, multiple seeds.

\textbf{Result.} Same-anchor 32/32; drift 0/32; explicit bridging 0/32.
Operations are anchor-bound: every reference frame needs its own
empirical evidence; a declared bridge does not transfer the operation.

\textbf{Token-level judgment.} The drift collapse is located at the tail
(``reads but cannot compute'': input head passes 32/32, the computation
positions collapse, 0.92 position accuracy at the tail) --- a distinct
mechanism fingerprint from E6's format-layer collapse.

\textbf{User intuition (07:13--07:43).} The user's direction steered the
whole design: ``the underlying choice is fine, the key is basepoint drift
--- what must be compared is basepoint drift; value\_zero should be the
basepoint's successor, right?'' The experiment compares exactly that:
drift vs same-anchor, with value-zero as the basepoint successor.

\textbf{Significance and declaration boundary.} This is the anchor-binding
law (law 1): operations bind to their basepoint; each reference frame
requires independent empirical evidence. The claim is about \emph{training
transfer under anchor change}, not about mathematical invariance --- the
mathematics of translation is trivially invariant, the training is not.
This asymmetry (math invariant, training bound) is the phenomenon the
framework formalizes.

\subsection{E12: succ structure matrix, two-pole (succ-matrix experiment)}\hfill\textit{(first success: 2026-08-15 08:44 bridging 32/32; reproduction 2026-08-16 14:54--15:04 per-seed identical)}

\textbf{Design reason.} E11 showed anchor binding. The same-basepoint
question remains: two succs at the same basepoint (succ$_a$, succ$_b$) ---
are they the same operation in different notations, or different
operations? The answer decides whether notation is translatable. The
two-pole design is deliberate: \emph{both} poles must be observed before
any claim (two-pole discipline I7ag) --- the zero-shot notation change
(0-pole) and the equivalence-declaration bridging (1-pole).

\textbf{Procedure.} Five groups $\times$ 3 seeds = 15 runs (dim 64, layers
2, epochs 60):
\begin{itemize}
\item[1a] baseline: train succ$_a$ (250 samples), OOD succ$_a$;
\item[1b] zero-shot: train succ$_a$, OOD \textbf{succ$_b$} (same basepoint,
notation change, no bridge);
\item[2] bridging: succ$_a$ + equivalence/anchor-error samples (270), OOD
\textbf{succ$_b$};
\item[2$'$] control: same 270 samples, OOD succ$_a$;
\item[4] symmetry: succ$_a$/succ$_b$ mixed (252), OOD mixed.
\end{itemize}

\textbf{Result.} All 45 runs (5 groups $\times$ 3 seeds $\times$ 3
independent runs) reach training accuracy 1.000. OOD: 1a 32/32; \textbf{1b
0/32}; \textbf{2 32/32}; 2$'$ 32/32; 4 32/32. The two-pole (0/32 vs 32/32)
never deviates across seeds or runs.

\textbf{Token-level judgment.} The 1b collapse is located at the
notation-branching position (the first succ$_b$ token); the input head
passes (32/32 input reconstruction), the failure is the notation fork.
Aggregate accuracy hides this; token-level reporting exposes the
mechanism fingerprint.

\textbf{User discipline applied.} Two-pole pre-declaration (I7ag): a
comparison may only be claimed when both poles are observed. 1b is the
0-pole, 2 is the 1-pole; the claim ``equivalence declarations bridge
notations'' is licensed only by the pair.

\textbf{Reproduction (2026-08-16 14:54).} The user asked to reproduce the
transformer experiments once more. E12 was re-run (15 runs), then run a
third time inside the self-contained package; all three runs agree
per-seed. The reproduction is part of the evidence chain, not a
re-derivation.

\textbf{Significance and declaration boundary.} Training intuition is
notation-bound statistics: 1.000 training accuracy does not transfer to an
untrained notation; OOD 1.0 proves only that statistics extrapolate within
that representation. Equivalence-declaration samples (pairwise
declarations; anchoring errors irrelevant) bridge the two notations ---
declarations make the information point at the correct structure. This is
the notation-translatability law (law 2) and the empirical seed of the
phase-locking claim (R138).

\subsection{E13: operation-distinction extrapolation (succ\_disc\_exp)}\hfill\textit{(first success: 2026-08-15 08:48, 56/56)}

\textbf{Design reason.} Can the training intuition distinguish operations
that share representation parts? The confusion pairs (2$\times$2 vs 2+2,
0$\times$2 vs 0+2) share every token except the operator. If operation
distinction is carried by the operator token alone, the distinction
extrapolates; if it depends on co-occurrence statistics, the OOD
confusion pairs collapse.

\textbf{Procedure.} Standard training (mixed operations, full 0 coverage);
OOD: confusion pairs (56 probes) on chain lengths/values inside the
training range but never co-occurring with the tested operator in the
training set.

\textbf{Result.} 56/56. Operation distinction extrapolates within the
representation: the operator token carries the distinction, and the
confusion pairs are resolved at every probe.

\textbf{Token-level judgment.} 56/56 full-sequence reconstruction; the
operator position and the result position both reconstruct correctly on
every confusion pair.

\textbf{Significance and declaration boundary.} This is the 1-pole of the
distinction question. Its 0-pole is E16 (no fake samples: 14/28, 50\%,
truth bias) --- the pair of experiments, not either one alone, establishes
the claim ``distinction requires the operator's negative evidence''.
Scoped: within-representation distinction extrapolates; cross-
representation distinction is a different question (E15).

\subsection{E14: symmetry center not at 0 (sym\_center\_exp)}\hfill\textit{(first success: 2026-08-15 06:00--06:16 midpoint OOD 1.000, epoch 13 (R163); formal re-run 2026-08-16, all OOD 100\% $\times$ 3 seeds)}

\textbf{Design reason.} The user's observation: ``some symmetries have
their axis not at 0 yet are natural numbers'' --- for a symmetric pair
$\{a,b\}$, the center $c=(a+b)/2$ is a natural number; the symmetry axis
(center) drifts from 0 to natural numbers. Does the training intuition
follow? This connects to R163 (symmetry as dimensionality reduction:
restoration point = midpoint) as a generation rule for the continuum
construction.

\textbf{Procedure.} Midpoint operation (midpoint token, chain-length
representation), three groups $\times$ 3 seeds: (1) full-range
extrapolation (train $a,b\in[0,8]$, OOD $[9,11]$); (2) center drift (train
$[0,6]$, OOD centers $\ge 5$, 72 probes); (3) symmetric-pair restoration
(train centers $\le 4$, OOD centers $5..7$).

\textbf{Result.} All OOD 100\% (10/10, 72/72, 24/24 $\times$ 3 seeds;
training accuracy 1.000). The symmetry axis (center) is \emph{not} an
anchor: the center is a computable result of chain-length arithmetic
$((\mathrm{len}(a)+\mathrm{len}(b))/2)$; with the structure fully expanded,
the center's drift from 0 to any natural number extrapolates.

\textbf{Token-level judgment.} Full-sequence reconstruction passes for all
probes; the midpoint token and the center numeral reconstruct
position-for-position in every group.

\textbf{Contrast with E11 (both poles).} Basepoint (chain bottom) = anchor
(drift collapses 0/32, E11); midpoint (symmetry axis) = generation rule
(drift stable). The two kinds of ``axis'' are essentially different --- a
distinction invisible to aggregate accuracy, exposed by the two-pole
comparison.

\textbf{Significance and declaration boundary.} R163's continuum
construction (dyadic midpoint generation) is empirically extrapolable in
the transformer. The claim is scoped to chain-length representation;
cross-representation transfer is tested separately (E15).

\subsection{E15: strict experiment (strict\_center\_exp)}\hfill\textit{(first success: 2026-08-16, zero-shot 26/26; cross-representation 0/26)}

\textbf{Design reason.} E14's honesty boundary: its group-2 OOD partially
overlapped trained $a,b$ values (combination extrapolation). The strict
version: fully zero-shot center drift (chain lengths never trained), plus
the cross-representation two-pole (I7ag: one 0-pole and one 1-pole before
claiming).

\textbf{Procedure.} Four groups $\times$ 3 seeds:
\begin{itemize}
\item[1] fully zero-shot center drift (train $a,b\in[0,4]$, OOD $[5,9]$,
chain lengths never trained);
\item[2] cross-representation stripping (train chain-length, OOD
position-value notation);
\item[3] cross-representation bridging (chain-length + numeric-equivalence
samples, OOD position-value);
\item[4] symmetric pairs fully zero-shot (train centers $\le 3$, arms
$\le 2$; OOD centers $4..6$, arms $3..4$).
\end{itemize}

\textbf{Result.}
\begin{center}
\begin{tabular}{llll}
\toprule
Group & OOD & token-level & judgment \\
\midrule
1 zero-shot center drift & 26/26 $\times$3 & 793/793 (1.00) & chain-length arithmetic zero-shot \\
2 cross-repr. stripping & 0/26 $\times$3 & 156/624 (0.25) & 0-pole: format layer (pos 4--14) \\
3 cross-repr. bridging & 15--25/26 & unstable & new finding: numeric equivalence $\neq$ operation equivalence \\
4 zero-shot symmetric pairs & 12/12 $\times$3 & --- & center+arms fresh, extrapolated \\
\bottomrule
\end{tabular}
\end{center}

\textbf{Token-level judgment.} Group 2's collapse is located at the format
layer (positions 4--14 all collapse; input head and read-out pairs pass)
--- the same fingerprint as E6's B/D collapse, confirming the format-layer
mechanism across experiments.

\textbf{New finding (group 3).} Numeric equivalence does not equal
operation equivalence: heterogeneous-representation bridging requires
operation-level equivalence samples. Same-isomorphism notations bridge
fully (E12); heterogeneous representations need operation-level
declarations (E6's isomorphism/heterogeneity boundary reappears at the
bridging dimension).

\textbf{Significance and declaration boundary.} The two-pole holds (1 and
4 at the 1-pole, 2 at the 0-pole): ``the symmetry axis extrapolates'' is
strictly scoped to the representation-internal intuition path ---
zero-shot extrapolation (chain-length arithmetic) succeeds, cross-
representation collapses (format layer), and equivalence samples bridge
only partially (operation-level equivalence required).

\subsection{E16: distinction-task 0-pole, no fake samples}\hfill\textit{(first success: 2026-08-16, 14/28 = 50\%)}

\textbf{Design reason.} E13 established the 1-pole of operation distinction
(56/56 under standard training). The coverage audit exposed a single-pole
problem: standard training includes negative evidence (fake results), so
56/56 could be carried by the negative samples alone. The 0-pole: train
\emph{without} fake samples (only true results), then OOD the confusion
pairs. If distinction requires negative evidence, the no-fake training
must collapse to chance.

\textbf{Procedure.} Training: true samples only (no fake results). OOD:
the same confusion pairs as E13 (28 probes).

\textbf{Result.} 14/28 = 50\% --- chance. Without conflicting information
the model cannot judge distinction (truth bias: the model defaults to
``true'' when it cannot distinguish).

\textbf{Token-level judgment.} The collapse is symmetric across the
confusion pairs: both directions (0$\times$2 vs 0+2 and their reverses)
fail at the operator position; no position-level asymmetry.

\textbf{Significance and declaration boundary.} Together with E13, the
two poles are complete: distinction requires the operator's negative
evidence; standard training provides it, no-fake training does not. The
claim is about \emph{evidence composition}, not about model capacity ---
the same model, same vocabulary, same sample count, differs only in the
presence of fake samples.

\subsection{E17: two-pole pre-declaration (methodology upgrade)}\hfill\textit{(first success: 2026-08-16, both poles complete)}

\textbf{Design reason.} E13/E14's single-pole problem was a process flaw:
the two-pole discipline (I7ag) was established \emph{after} those
experiments were run, so their 0-poles had to be added retroactively
(E16). E17 upgrades the methodology: the 0-pole and the 1-pole are
\emph{pre-declared at design time}, before any training runs.

\textbf{Procedure.} Three groups, each with a pre-declared 0-pole and
1-pole (six conditions), $\times$ 3 seeds:
\begin{itemize}
\item[a] equivalence: chain-length vs bare-symbol representation of
equivalence (1-pole: chain-length; 0-pole: bare symbol);
\item[b] distinction: true+fake vs no-fake training on the confusion pairs
(1-pole: 28/28; 0-pole: 14/28);
\item[c] center: chain-length midpoint vs cross-representation midpoint
(1-pole: 10/10; 0-pole: 0/26).
\end{itemize}

\textbf{Result.} Every group shows the full two-pole column: a 0-pole
collapse and a 1-pole pass, both pre-declared. (a) 50/50 vs 2/50; (b)
28/28 vs 14/28; (c) 10/10 vs 0/26.

\textbf{Token-level judgment.} Each pole reported token-level (sample
pairs + position-level collapse localization) per the reporting
discipline; no aggregate-only claims.

\textbf{Significance and declaration boundary.} The methodology law: a
comparison is claimed only when both poles were declared before the runs
and both were observed after. The upgrade is itself part of the trail ---
the discipline's history (E13/E14 retroactive, E16/E17 pre-declared) is
recorded, not erased.

\subsection{E18: mechanics-token full set}\hfill\textit{(first success: 2026-08-16, paired 5/9 vs one-way 1/9)}

\textbf{Design reason.} The framework claims directions must be declared in
pairs (R136). The experiment tests this at the token level: 8 mechanics
tokens (succ/add/sub/mul/div/mid/recip/inv) in the same chain-length
representation (single structure, not mixed --- the E17a representation
lesson). Two training regimes: paired declaration (inverse pairs trained
together) vs one-way (only the forward direction). OOD: new-range
confusion pairs over all tokens.

\textbf{Procedure.} 1-pole: paired inverse training; 0-pole: one-way
training. OOD identical (new-range confusion pairs, all-token pairs).

\textbf{Result.} Inverse-knowledge discrimination 5/9 (paired) vs 1/9
(one-way): paired declaration makes difference decidable, one-way
collapses --- R136's empirical witness. Forward extrapolation passes in
both regimes (8/8, chain-length statistics). recip/inv fail in both
regimes (4 probes): 2 samples are too few (EXP-SYM-INV replication) plus
structural poverty of the natural-number domain.

\textbf{Token-level judgment.} The paired advantage is located at the
inverse-pair positions (the discriminator needs both directions to
separate the pair); forward positions pass in both regimes.

\textbf{Significance and declaration boundary.} Direction declarations
are empirically effective in pairs; single declarations leave the
discrimination undecidable. The claim is scoped to the 8 tested tokens and
the chain-length representation; recip/inv's failure is a data-diversity
statement (2 samples), not a token-design claim.

\section{Five structural laws}

The E-series converges on five laws, each with its two-pole evidence:

\begin{enumerate}
\item \textbf{Anchor binding (E11, first success 2026-08-15 07:43).}
Operations bind to their basepoint; reference frames require independent
evidence (same-anchor 32/32, drift 0/32, bridging 0/32). Math is
invariant, training is bound --- the asymmetry the framework formalizes.
\item \textbf{Notation translatability (E12, first success 2026-08-15
08:44).} Same-basepoint notation change collapses zero-shot (0/32) and
bridges with equivalence declarations (32/32); notation is translation
only via declarations. (R138 phase locking.)
\item \textbf{Shift invariance (E9, first success 2026-08-15 07:33).} The
successor rule is translation-invariant within an axis (succ(0) alone
yields succ(1..7) 14/14); extrapolation is chain-length arithmetic.
\item \textbf{Structure-bearing representation (E6, first success
2026-08-15 08:36--08:37).} Position-value collapses (0/30), true numerals
extrapolate (30/30); structure must be carried by the representation, not
assumed (E7/E8: 0 = chain-length 0, zero-shot 45/45).
\item \textbf{Definition-layer invisibility (E4/E5, first success
2026-08-15 06:54).} The definition of 0 and the design of succ are
invisible to training (24/24 runs equivalent); training reads
representation form, not definitions.
\end{enumerate}

Each law is stated with its polarity evidence and its scope; no law is
claimed beyond its representation (chain-length, dim 64, layers 2).

\section{From laws to formalized framework}

The laws map to the Pat framework (basepoint-relative operator semantics):

\begin{center}
\begin{tabular}{lll}
\toprule
Law & Framework claim & Lean file (0 sorry) \\
\midrule
Anchor binding & R136 directions declared in pairs & \\
Notation translatability & R138 phase locking & MutualLocking.lean (R143) \\
Equivalence declaration & R143 interlock matrix det $\neq$ 0 & MutualLocking.lean \\
Basepoint = anchor & basepoint-relative operators & BasepointMove.lean (C026) \\
Continuum & R150/R151/R163 & PatCountableInfinitPhaseUnification.lean \\
 & & ContinuumPatGrid.lean, ContinuumConstructible.lean \\
\bottomrule
\end{tabular}
\end{center}

The formalization: claims R136--R153 (the framework main chain) all PROVED
with no \texttt{sorry}, completed 2026-08-13 13:31 (18 claims); the Toolkit
branch extends to 150+ Lean files, claims R001--R232 (with 43 Ruler*
experience-layer claims and TK1--TK5); full \texttt{lake build} passes
(3631 jobs, mathlib v4.32.2, lean-toolchain pinned). Key theorems:
$J^2 = -I$ (symplectic rotation, R149 four-phase interlock, $S^3$),
the mechanics-symmetry refinement R233 (2026-08-16, 29 theorems, 0
\texttt{sorry}): the four layers of mechanics symmetry, each interlocked
symmetry pair orthogonal, the symplectic structure ($J^2=-I$; each degree
of freedom = a four-phase interlock pair), the full symplectic embedding
(12 pairs $= \mathbb C^{12}$), and the gravitational symmetry direction
(structure layer has no privilege $\wedge$ action layer is one-way),
mutual locking $\iff$ determinant $\theta(r+1/r) \neq 0$ (R143), the
countable-infinity/continuum phase-locking unification (R150), the continuum
as Pat-grid closure (R151), and the endogenous continuum construction with
zero Nat in the definition layer (R163). The ZeroRelative branch formalizes
the basepoint-relative algebra (C001--C025: heap retract, semiconjugacy,
displacement equivalence, endogenous generation) and the operator semantics
(C026--C038: every operator has its own basepoint, $\pm1$ as projection
extremes, symplectic projection, Q2\_0 grid as sign encoding, state-matrix
complex structure).

\section{The logical completeness theorem and the three mapping-judgment theorems}

\subsection{The logical completeness theorem (2026-08-15 09:14--09:47)}

At 09:14 the user asked to review the whole process against the framework's
logical completeness theorem and to design a formal verification project:
``how do you run a logically complete experiment? And I conjecture that
orthogonality itself is a symmetry --- can a single axis complete training?
That needs to be compared.'' The theorem, as extracted from the claims
ledger (R031, completeness C1--C7), states:

\begin{quote}
\emph{The decoupling operator $D$ (R1 exclusion + R2 cancellation + R3
layering) $\Longrightarrow$ completeness at any order, any element, any
subject-object: perfect OOD.}
\end{quote}

The verification project (same day):
\begin{itemize}
\item 09:16 --- orthogonality-declaration bridging succeeds same-anchor
  (32/32); cross-anchor orthogonality is seed-dependent (B/C: 32/32, 32/32,
  0/32), in contrast to the same-anchor stability of E12;
\item 09:17 --- the \emph{double} declaration (orthogonality + equivalence)
  restores stability: 5/5 across seeds. A single declaration is insufficient;
\item 09:18---09:19 --- the user drew the corollary (I7u): if orthogonality
  is a symmetry, there is no axis orthogonal to every axis; an ``orthogonal
  axis'' is a phase on an orthogonality symmetry axis. Then the stronger
  claim (I7v): \emph{two groups of symmetries are necessary to localize a
  direction} --- Cartan--Dieudonn\'e: rotations are composites of two
  reflections, a single reflection only defines $\pm$; therefore Pat0's
  declarations must be two groups of symmetries, and any thought-experiment
  based on a single dimension is invalid. Direction localization with two
  symmetry groups: OOD 16/16 $\times$ 3 seeds (09:19);
\item 09:22 --- the user asked for the minimum number of symmetry groups when
  the declared directions are not two orthogonal axes (conjecture: 3 or 4,
  because an orthogonal phase implies an angle mapping);
\item 09:30--09:36 --- the completeness dimensions (order / element /
  subject-object) re-verified token-level per the new reporting discipline
  (algebraic axioms 12/12, involutions 12/12, truth/false separation 6/6);
\item 09:19 --- the methodology was written up as the short paper
  \emph{Logical-Completeness Experimental Design Methodology}, later merged
  into the release package (``self-redemption series'').
\end{itemize}

\subsection{The three mapping-judgment theorems (2026-08-16 16:16--16:23)}

On 2026-08-16 the judgment theorems for the token system were located as
\emph{multi-to-multi mapping judgments}. The user's correction at 16:17
fixed the two-theorem positioning:
``the symbol-norm theorem is: how does one judge the multi-to-multi mapping
from symbols to logic and intuition; the presentation-legality theorem is:
how does one judge the multi-to-multi mapping from logic and intuition to
presentation.'' The three theorems:

\begin{enumerate}
\item \textbf{Symbol-norm theorem} (symbols $\to$ logic/intuition): every
  mapping target must be resolvable --- ambiguity converges by context
  (no context = all candidates; given context = one); unresolvable
  ambiguity = non-norm. Two symmetric directions (A: symbol $\to$ logic;
  B: symbol $\to$ intuition).
\item \textbf{Presentation-legality theorem} (logic/intuition $\to$
  presentation): every presentation target must be resolvable (notation
  context: precedence/associativity, concrete notation in the P layer);
  numeric coincidence without declaration is illegal (equivalence bridging
  must be declared). Two symmetric directions.
\item \textbf{Intuition-precision theorem} (intuition $\to$ logic/
  presentation): the intuition-token interface must be definitionally
  resolvable (does not impersonate a logic concept); intuition-path
  precision is judged by two-pole comparison; polarity-confusion events are
  the canonical failure.
\end{enumerate}

A fourth theorem (symbol-presentation mapping) was added at 16:19--16:20,
completing the five-layer mapping system (B/C/S/G/P). All four were
formalized in \texttt{MappingJudgmentTheorems.lean} (5 theorems, 0
\texttt{sorry}, \texttt{lake build} 2997 jobs) with an experiment-coverage
matrix (E6/E4/E5/E12/E13/E15/E14 + R065), committed 16:23.

\section{Timestamped evidence chain}

The trail is anchored by six independent evidence sources (all in the
artifact):

\begin{enumerate}
\item Session records (db.sqlite message/part, per-message timestamps,
UTC+8; 55{,}047 messages, 2026-07-15 onward) --- the primary source;
every timestamp in this paper is traceable to a session record.
\item Task execution log: 2303 turns (turn\_usage; later extended to
2773).
\item Model-viewpoint phase records (daily) --- the API model's own
claims separated from the user's, per the distillation-pressure
corrections.
\item git log commit timestamps (first Lean commit 2bc157e, 2026-08-07
10:53; first Lean no-\texttt{sorry} C001 at 11:31 the same day;
Riemann-direction archive ec87658 2026-08-12 02:54; full-project track
df65309 2026-08-13 17:48).
\item Experiment artifact file timestamps (E1--E13 results, 2026-08-15
07:22--09:00; result files in the self-contained package, cross-checked
against session records).
\item Lean no-\texttt{sorry} timeline (R136--R153 completion 2026-08-13
13:31; 786 no-\texttt{sorry} records extracted and deduplicated on
2026-08-16).
\item The four-phase classification record (2026-08-16): all 4068 user
messages from 2026-07-29 classified as positive / leaning-positive /
neutral / leaning-negative / negative (13 / 101 / 2159 / 37 / 1758 after
deduplication to 2158) --- an independent audit of which statements were
the user's affirmations and which were the model's.
\end{enumerate}

Key anchors (each with its session record): 2026-07-29 22:39:12 (UTC+8)
first \emph{imperfect} generalization (halving experiment, 450 samples);
2026-07-29 22:52 strongest early generalization (train 5 relations leave
1: $\le$ 81\% $\to$ 90\%); 2026-08-06 17:10--18:46 Riemann visualization
(delivery 17:55); 2026-08-07 10:53/11:31 Lean first commit / first
no-\texttt{sorry}; 2026-08-12 00:00--02:54 Riemann formalization session
(C011--C025); 2026-08-12 framework core claims (R134--R153) and the
Riemann publication (09:43--11:56); 2026-08-13 13:31 all 18 framework
claims PROVED; 2026-08-13 14:19--18:37 exercise series start and first
publication (10 DOIs); 2026-08-14 12:47 exercise re-qualification;
2026-08-15 04:01 first \emph{perfect} generalization; 2026-08-15
07:22--09:00 E1--E13; 2026-08-15 09:14--09:47 logical-completeness
verification project; 2026-08-16 14:54--15:04 E12 reproduction;
2026-08-16 16:16--16:23 three mapping-judgment theorems; 2026-08-16
self-contained package.

\section{Human corrections steering the API execution}

The trail is also a record of how the user's instructions corrected the API
model's execution. We list the documented corrections (each with its session
timestamp) because they are part of the evidence: they show which statements
are the user's claims and which are the model's insistence.

\begin{itemize}
\item \textbf{2026-08-12 09:44} --- referent correction: ``no, the main paper
of the Riemann formalization''. The model had picked the wrong document
(Unified Framework instead of the Riemann main paper); the user
disambiguated; the DOI was added to the right paper.
\item \textbf{2026-08-12 11:08} --- attribution correction: the user's
instruction that the Riemann paper must state explicitly: no claim of
proving the Riemann hypothesis / all results are restatements of known facts
/ novelty = KNOWN. (Later reinforced at 2026-08-13 03:05: ``revise the
Riemann paper wording: every claim about `proving' is what DeepSeek insisted
on'' --- the unproved statements are the model's insistence, not the user's
claim.)
\item \textbf{2026-08-13 14:24} --- positioning correction (exercise I): ``not
the computation chapter of the psyche paper, it is a homework exercise for
the foundation chapter''.
\item \textbf{2026-08-13 14:31} --- scope correction (exercise I): ``do not
over-enumerate, meaningless; under this framework we trade speed for
fuzziness and recover precision with structure'' --- the 9-theorem
enumeration was deleted, leaving 4 mechanism theorems.
\item \textbf{2026-08-13 14:48} --- citation instruction (exercise I): ``cite
as many living mathematicians, philosophers, linguists and scientists as
possible'', with every citation verified (no citations from memory).
\item \textbf{2026-08-14 12:47} --- series re-qualification: all exercises
must carry ``exercise / topic / Lean / DOI'' and declare they represent no
mathematical conclusion, only observation reports based on Lean
formalization.
\item \textbf{2026-08-15 09:14} --- completeness-check instruction: review the
whole process against the logical completeness theorem and design a formal
verification project; conjecture that orthogonality is itself a symmetry
(single-axis training must be compared). Generated the 09:16--09:19
double-declaration result and the two-symmetry-groups corollary (I7v).
\item \textbf{2026-08-15} --- two-pole discipline (I7ag): a comparison may
only be claimed when \emph{both} poles are observed (one 0-OOD and one
1-OOD). E13's single-pole design was flagged by this rule and re-done as
E16/E17.
\item \textbf{2026-08-15} --- token-level reporting discipline: aggregate
accuracy (0.99/1.0) is not sufficient; every claim must report full sample
logs, token pair/total counts, and position-level accuracy (collapse
localization). token\_report.py was built under this rule.
\item \textbf{2026-08-15} --- intuition/structure stripping discipline:
before a complete ablation comparison, every OOD result must be assumed to be
an intuition-path result; structural stripping is mandatory.
\item \textbf{2026-08-16 16:17} --- theorem-positioning correction: ``the
symbol-norm theorem is the judgment of the multi-to-multi mapping from
symbols to logic and intuition; the presentation-legality theorem is the
judgment of the multi-to-multi mapping from logic and intuition to
presentation''. The three mapping-judgment theorems were re-located under
this definition.
\item \textbf{2026-08-16 06:17} --- record-format correction: ``you need to
improve this: the time the motivation was proposed, the time the experiment
was completed (first completion counts, e.g., the basepoint drift in the
Riemann direction), and the time the paper was written''. This rule
generated the per-experiment first-success timestamps used throughout this
paper.
\item \textbf{2026-08-16} --- timeline correction: the first generalization
is 2026-07-29 (22:39:12 UTC+8); the model had misdated it; the user
corrected the record.
\item \textbf{2026-08-16 06:18 / 06:21} --- attribution correction under
distillation pressure: ``you have already claimed it as your own, you
know?'' and ``it is not just this: you have already published what we
discussed''. These corrections fixed the record's authorship boundary.
\item \textbf{2026-08-16 14:54} --- reproduction instruction: ``reproduce our
transformer experiments once more'' --- E12 re-run (15 runs, per-seed
identical) + third in-package run.
\end{itemize}

\section{The user's questions, corrections, and examples}

The trail is driven by the user's questions. This section records them
verbatim (with session timestamps), because they are the actual research
hypotheses --- each question determined the next experiment, and each
correction fixed the record. The API model's role was execution and
mechanical verification; the user's role was direction, judgment, and
example.

\subsection{Questions that steered the trail (verbatim, from the session log)}

\begin{itemize}
\item \textbf{2026-07-29 18:46} --- the first experiment: ``we train an
addition model, without providing the definition of addition, only the
computation cases within 1000, and see whether 5-digit addition can be
learned'' (the pure next-token baseline: 2-digit 100\%, 3-digit 0\%).
\item \textbf{2026-07-29 22:20} --- the information bound: ``the sample
count should be calculated by the minimum guessing-game steps'' (Mastermind
mechanism; led to the information-gain sampling and the halving
experiment).
\item \textbf{2026-07-29 22:39} --- the first generalization confirmation:
``this means generalization happened, the necessary number of steps was not
needed to get the necessary result. What is your halving logic, and which
logic tokens generalized and which failed?'' (the first per-token
analysis, 22:41).
\item \textbf{2026-07-30 00:00} --- the research roadmap: ``in math and
logic, propositions are consistent (no self-reference, no contradiction);
once we enter medicine, social science, law, there are many contradictory
propositions that must live in one framework. This needs research --- but
first we solidify math and logic, by code practice.''
\item \textbf{2026-07-30 00:42} --- the base-change conjecture: ``I suspect
that base-change equivalence naturally carries multi-digit out'' (base
change $\to$ addition generalization 81\%, multiplication 12\% --- the
latter requiring its own training, as the user then corrected: ``of course
multiplication must be trained'').
\item \textbf{2026-07-30 05:30} --- the parallel-token principle: ``note it
is parallel \emph{logic tokens}, not parallel elements --- once all logic
tokens of an element are sufficiently trained through parallel logic
tokens, the element generalizes even if never trained'' (verified: leaving
$\le$ with a broken counterpart axis 55\% vs 90\% with complete
counterparts).
\item \textbf{2026-08-12 00:14} --- the honesty question: ``wait, so along
this line of thought can the Riemann hypothesis be proved?'' (the model
answered: no; all six theorems are OBSERVATION/KNOWN restatements, none
touches the zeta function or zeros --- the honest boundary is part of the
record).
\item \textbf{2026-08-12 00:16} --- the true-complex-axis demand: ``we
cannot construct the prime axis, but we can construct the complex axis
from lambda calculus, right? Note I want the \emph{true} complex axis, not
the fake one where each successor is the natural 1'' (the user rejects
axes whose successor is pre-assumed to be the natural numbers; the true
axis is multiplicative iteration --- rotation-scaling orbits).
\item \textbf{2026-08-12 00:17} --- the projection hypothesis: ``for the
construction of roots: suppose the complex axis is essentially the
projection of the higher-dimensional structure $-1$ into human
mathematical space; under this assumption, can it be constructed?''
(formalized as the $a+bJ$, $J^2=-1$ rotation algebra, 00:19--00:21).
\item \textbf{2026-08-12 00:23} --- the basepoint-drift statement: ``since
our constructed complex plane carries the naturals on its axis --- do you
know what this means? It means our lambda basepoint has moved onto the
complex axis, i.e., the basepoint has drifted; its position is not the
origin but some transform of $i$''.
\item \textbf{2026-08-12 00:25} --- the fake-real-axis statement: ``then
that so-called real axis is very suspicious; I suspect it is not a real
number axis at all, but another fake complex axis, right?''.
\item \textbf{2026-08-15 06:54} --- the numeral intuition: ``numeral is
itself intuition; the true numeral is the successor chain from the
basepoint ... and 0-multiplication must be compared against nonzero
multiplication'' (E8's zero-shot design).
\item \textbf{2026-08-15 07:13} --- the drift comparison: ``the underlying
choice is fine; the key is basepoint drift --- what must be compared is
basepoint drift. value\_zero should be the basepoint's successor, right?''
(E11's design).
\item \textbf{2026-08-15 08:36} --- the stripping discipline: ``before a
complete ablation comparison, every OOD result must be assumed to be
intuition; differences exist only through comparison''.
\item \textbf{2026-08-15 09:14} --- the orthogonality conjecture: ``review
the whole process against our logical completeness theorem, design a
formal verification project on how to run logically complete experiments.
Then a conjecture: I suspect orthogonality itself is a symmetry --- can a
single axis complete training? It must be compared.''
\item \textbf{2026-08-15 09:19} --- the two-symmetry-groups corollary:
``orthogonality as symmetry yields the conclusion: two groups of
symmetries are necessary to localize a direction; therefore Pat0's
declaration must be two groups of symmetries; any thought-experiment based
on a single dimension is wrong'' (verified: direction localization 16/16
$\times$ 3 seeds).
\item \textbf{2026-08-15 09:22} --- the minimum-groups question: ``if the
declared directions are not two orthogonal axes, how many groups of
symmetries are needed at minimum? I suspect 3 or 4, because if an
orthogonal symmetry has phase, then angle has a mapping''.
\item \textbf{2026-08-16 14:54} --- the reproduction instruction:
``reproduce our transformer experiments once more''.
\end{itemize}

\subsection{Examples citing the mathematicians' consensus}

The user's examples often cite established mathematical consensus and then
re-interpret it under the projection hypothesis --- the pattern of the
whole trail:

\begin{itemize}
\item \textbf{2026-08-07 (thinking document), translation vs inversion}:
``mathematicians have told me that 1 and $1/a+1/b$ are surely symmetric,
translation and inversion; I believe their construction. But under my
hypothesis, $1/a+1/b$ is the projection of inversion-1 after the
projection lost structure. If my hypothesis holds, it can only mean one
thing: the basepoint has drifted --- the lambda-calculus origin has moved
around the circle of radius 1 centered at 1''. The consensus (the
symmetric pair) is preserved; the \emph{mechanism} (why they are
symmetric) is re-derived from projection loss.
\item \textbf{2026-08-07 (thinking document), irrationals}: ``I firmly
believe irrationals are the remainder of some higher-dimensional structure
after projection into our space lost part of the information-carrying
structure. The reason is simple: every intuition-path construction moves
along a direction away from the basepoint; only irrationals need to walk
back. The unit ball in the higher dimension is complete 1 in one
direction; projected into our world it leaves only $\sqrt2/2$ --- this is
likely the irrational's 1''.
\item \textbf{2026-08-09 01:27 (ChatGPT conversation), Church numerals}:
``the lambda-calculus basepoint self-application iteration presupposes a
concept distinct from the basepoint: the self-application operation,
$\mathrm{lambda}(\mathrm{base},\mathrm{base})$. But in the concept I want
to experiment with, there is no such operation: $\mathrm{base}(\mathrm
{base},\mathrm{base})$'' --- the Church-numeral consensus is rejected as
impure: it smuggles an implicit operation role, while the pure-relative
design occupies operator/operand/basepoint with one primitive.
\item \textbf{2026-08-15 09:19, Cartan--Dieudonn\'e}: the user's
two-groups corollary was checked against the classical theorem that
rotations are composites of two reflections (a single reflection defines
only $\pm$); the consensus confirms the corollary rather than the
framework being derived from it --- the direction-localization experiment
(16/16 $\times$ 3) is the witness.
\end{itemize}

\subsection{Corrections during experiments (verbatim)}

Beyond the record-format and attribution corrections of
Section~9 (Human corrections), the experiment log contains execution
corrections that directly changed designs:

\begin{itemize}
\item \textbf{2026-07-29 22:42} --- output-format correction: ``my
requirement was to output the \emph{element's logic tokens}, not the
element! And this means our prefix question format needs more space'' ---
the model had been outputting truth values; the four-question-type system
followed (22:49).
\item \textbf{2026-07-29 22:54} --- architecture correction: ``the data
synthesizer, result translator, and result analyzer must all be
pluggable; changing one position in the concept token takes effect
globally'' --- the JSONL-driven CLI architecture.
\item \textbf{2026-07-30 08:00} --- ``legacy's logic tokens are incomplete;
the information about `='' was not trained/sampled --- do you think this
needs a logic token or training samples?'' (both, answered by the
experiment: 47\% $\to$ 79\% with completed counterparts).
\item \textbf{2026-08-15 06:15} --- ``a single operation with small
samples cannot be learned anyway'' --- the failed single-operation run
(OOD 0/15) was corrected to mixed training before E3's 36/36.
\item \textbf{2026-08-16 16:17} --- the theorem-positioning correction for
the three mapping-judgment theorems (Section~8.2).
\end{itemize}

\section{Reproducibility and provenance}

The artifact is a self-contained package (also published as a standalone
repository):
\begin{itemize}
\item Training: E12 one-command reproduction (5 groups $\times$ 3 seeds,
GPU auto, 5--10 minutes); E17 two-pole pre-declaration; E15 strict
experiment; full E1--E13 (~30 minutes). All results re-generate the
package's report tables.
\item Formalization: \texttt{lake build} (pinned lean-toolchain
v4.32.2 + mathlib v4.32.2), 0 \texttt{sorry}.
\item E12 reproduction fidelity: three independent runs (original $\times$2 +
in-package) per-seed identical, no random drift.
\item Provenance methodology: SHA-256-anchored release manifests, per-claim
SHA256 publication, per-claim anonymous mirror (dual-repo: real +
anonymized), and a publish pipeline (scan $\to$ SHA-256 digest $\to$ manifest $\to$
real repo $\to$ anonymize $\to$ blind repo).
\end{itemize}

\section{Related work and conclusion}

The empirical part relates to interpretability and OOD generalization studies
of small transformers (known: OOD extrapolation depends on representation;
cf.~\cite{power2022grokking,vaswani2017attention}); the formal part builds on
heap/torsor algebra \cite{bergman2009universal} and the previously submitted
Riemann-direction formalization (C011--C025). The trail as a whole ---
observation $\to$ laws $\to$ framework $\to$ machine-checked theorems, with
timestamped evidence --- is the contribution.

\paragraph{Conclusion.} We record a complete research trail: from the 0/1 OOD
dichotomy of a low-epoch transformer (2026-07-29, 22:39:12; E12: 0/32 vs
32/32 across 45 runs) through eighteen individually documented experiments,
five structural laws, and the Pat framework, to a Lean formalization
(R136--R153 all PROVED, no \texttt{sorry}; 3631 jobs; R150/R151/R163
continuum chain). Three parallel threads are part of the same trail: the
Riemann direction (2026-08-06 visualization $\to$ 2026-08-12 C011--C025
formalization and publication), the homework-exercise series (2026-08-13
$\to$ 08-15, 25 reports, all Lean 0-\texttt{sorry}), and the verification
methodology itself --- the logical-completeness project (2026-08-15
09:14--09:47: decoupling operator $D$ = R1 exclusion + R2 cancellation +
R3 layering, completeness at any order/element/subject-object, the
two-symmetry-groups corollary) and the three mapping-judgment theorems
(2026-08-16 16:16--16:23, \texttt{MappingJudgmentTheorems.lean} 0
\texttt{sorry}).

The trail is timestamped by seven independent evidence sources and
reproducible from a self-contained package. Its writing follows the
recording disciplines of the trail itself: every passage states its
premises before its conclusion (no logical leaps: a claim is licensed only
by its two-pole evidence); every symbol, number, and timestamp is
traceable to a session record (no unanchored notation); every statement
carries its provenance --- the user's instruction, the model's insistence,
or the mechanical verification --- (no unverifiable assertion); and the
user's intuitions are recorded separately from the API model's
interpretations, with the corrections that separated them (no conflation
of intuition path and logic path). Under provenance pressure, we argue,
recording the \emph{path} --- not only the destination --- is the
defensible form of priority: each experiment, each claim, and each theorem
carries its own timestamped anchor, and each anchor is either a session
record, a commit, a file timestamp, or a machine-checked proof.

\section*{Theorem inventory (selection)}

\begin{center}
\small
\begin{tabular}{lll}
\toprule
Claim & File & Content \\
\midrule
R136 & (framework) & directions declared in pairs $\{d,-d\}$ \\
R138 & PhaseRelationLocking.lean & phase relation locking \\
R143 & MutualLocking.lean & interlock $\iff$ det $\neq$ 0 \\
R149 & Pat4Phase.lean & four-phase interlock, $S^3$, $J^2=-I$ \\
R150 & PatCountableInfinitPhaseUnification.lean & countable/continuum \\
R151 & ContinuumPatGrid.lean & continuum = Pat-grid closure \\
R163 & ContinuumConstructible.lean & endogenous, zero Nat \\
C001--C010 & ZeroRelative/ & heap, semiconjugacy, generation \\
C026--C038 & ZeroRelative/ & operator semantics, Q2\_0, state structure \\
\bottomrule
\end{tabular}
\end{center}

\bibliographystyle{ACM-Reference-Format}
\bibliography{refs}

\end{document}
